\documentclass[11pt,letterpaper]{article}
\usepackage[utf8]{inputenc}
\usepackage[spanish, es-tabla]{babel}
\usepackage[version=3]{mhchem}
\usepackage[journal=jacs]{chemstyle}
\usepackage{amsmath}
\usepackage{amsfonts}
\usepackage{amssymb}
\usepackage{makeidx}
\usepackage{xcolor}
\usepackage[stable]{footmisc}
\usepackage[section]{placeins}
%Paquetes necesarios para tablas
\usepackage{longtable}
\usepackage{array}
\usepackage{xtab}
\usepackage{multirow}
\usepackage{colortab}
%Paquete para el manejo de las unidades
\usepackage{siunitx}
\sisetup{mode=text, output-decimal-marker = {,}, per-mode = symbol, qualifier-mode = phrase, qualifier-phrase = { de }, list-units = brackets, range-units = brackets, range-phrase = --}
\DeclareSIUnit[number-unit-product = \;] \atmosphere{atm}
\DeclareSIUnit[number-unit-product = \;] \pound{lb}
\DeclareSIUnit[number-unit-product = \;] \inch{"}
\DeclareSIUnit[number-unit-product = \;] \foot{ft}
\DeclareSIUnit[number-unit-product = \;] \yard{yd}
\DeclareSIUnit[number-unit-product = \;] \mile{mi}
\DeclareSIUnit[number-unit-product = \;] \pint{pt}
\DeclareSIUnit[number-unit-product = \;] \quart{qt}
\DeclareSIUnit[number-unit-product = \;] \flounce{fl-oz}
\DeclareSIUnit[number-unit-product = \;] \ounce{oz}
\DeclareSIUnit[number-unit-product = \;] \degreeFahrenheit{\SIUnitSymbolDegree F}
\DeclareSIUnit[number-unit-product = \;] \degreeRankine{\SIUnitSymbolDegree R}
\DeclareSIUnit[number-unit-product = \;] \usgallon{galón}
\DeclareSIUnit[number-unit-product = \;] \uma{uma}
\DeclareSIUnit[number-unit-product = \;] \ppm{ppm}
\DeclareSIUnit[number-unit-product = \;] \eqg{eq-g}
\DeclareSIUnit[number-unit-product = \;] \normal{\eqg\per\liter\of{solución}}
\DeclareSIUnit[number-unit-product = \;] \molal{\mole\per\kilo\gram\of{solvente}}
\usepackage{cancel}
%Paquetes necesarios para imágenes, pies de página, etc.
\usepackage{graphicx}
\usepackage{lmodern}
\usepackage{fancyhdr}
\usepackage[left=4cm,right=2cm,top=3cm,bottom=3cm]{geometry}

%Instrucción para evitar la indentación
%\setlength\parindent{0pt}
%Paquete para incluir la bibliografía
\usepackage[backend=bibtex,style=chem-acs,biblabel=dot]{biblatex}
\addbibresource{references.bib}

%Formato del título de las secciones

\usepackage{titlesec}
\usepackage{enumitem}
\titleformat*{\section}{\bfseries\large}
\titleformat*{\subsection}{\bfseries\normalsize}

%Creación del ambiente anexos
\usepackage{float}
\floatstyle{plaintop}
\newfloat{anexo}{thp}{anx}
\floatname{anexo}{Anexo}
\restylefloat{anexo}
\restylefloat{figure}

%Modificación del formato de los captions
\usepackage[margin=10pt,labelfont=bf]{caption}

%Paquete para incluir comentarios
\usepackage{todonotes}

%Paquete para incluir hipervínculos
\usepackage[colorlinks=true, 
            linkcolor = blue,
            urlcolor  = blue,
            citecolor = black,
            anchorcolor = blue]{hyperref}

%%%%%%%%%%%%%%%%%%%%%%
%Inicio del documento%
%%%%%%%%%%%%%%%%%%%%%%

\begin{document}
\renewcommand{\labelitemi}{$\checkmark$}

\renewcommand{\CancelColor}{\color{red}}

\newcolumntype{L}[1]{>{\raggedright\let\newline\\\arraybackslash}m{#1}}

\newcolumntype{C}[1]{>{\centering\let\newline\\\arraybackslash}m{#1}}

\newcolumntype{R}[1]{>{\raggedleft\let\newline\\\arraybackslash}m{#1}}

\begin{center}

%Reemplace el texto en cada uno de los siguientes espacios. No borre los corchetes porque se producen errores en la compilación del documento.

\textbf{\LARGE{Título de la Práctica}}\\
	\vspace{7mm}
		\textbf{\large{Primer autor}}\\
	\textbf{\large{Segundo autor}}\\
	\vspace{4mm}
	\textbf{\large{Grupo A1 o G1}}\\
	\today
\end{center}

\vspace{7mm}

\section{Introducción}

Esta sección debe incluir información general que explique la importancia del experimento. Se deben utilizar referencias para soportar la base científica del mismo; esto le da credibilidad, porque indica que preparó adecuadamente la práctica. Debe incluir la información relevante relacionada con los experimentos; haga especial énfasis en condiciones inusuales o críticas. Las referencias deben ir numeradas consecutivamente en el texto y aparecer como una lista en la sección~\nameref{sec:references}.\autocite{Cooper2009} Las referencias más comunes incluyen libros, manuales de laboratorio y artículos.

\section{Resultados}

%En esta sección se debe presentar de forma condensada los datos obtenidos durante el experimento. Todos los datos deben ser resumidos en tablas y/o gráficas enumeradas. El título debe contener una descripción clara del contenido. 
%Las tablas deben reservarse para reportar listas de valores, mientras que la relaciones entre variables deben reportarse en gráficas. Tanto las gráficas, los esquemas de reacción, las fotografías y las ilustraciones deben agruparse bajo el título figura.

\begin{longtable}{C{3cm}C{4cm}C{4cm}C{4cm}}
\caption{Puntos de ebullición de las muestras problema.}
\label{tab:boiling}\\
\hline
\bfseries{Compuesto} & \bfseries{Punto de ebullición experimental \SI{}{\degreeCelsius}} & \bfseries{Punto de ebullición corregido \SI{}{\degreeCelsius}} & \bfseries{Punto de ebullición literatura \SI{}{\degreeCelsius}}\\
\hline
\endfirsthead
\multicolumn{2}{c}{\bfseries \tablename{} \thetable \ldots continuación}\\
\hline
\bfseries{Compuesto} & \bfseries{Punto de ebullición experimental \SI{}{\degreeCelsius}} & \bfseries{Punto de ebullición corregido \SI{}{\degreeCelsius}} & \bfseries{Punto de ebullición literatura \SI{}{\degreeCelsius}}\\
\hline
\endhead
\hline
\multicolumn{2}{r}{\ldots continúa en la siguiente página}\\
\endfoot
\hline
\endlastfoot

\includegraphics{./images/toluene.eps} & b & c & \num{110.6} \\ 
\includegraphics{./images/cyclohexane.eps} & b & c & \num{80.74} \\ 
\includegraphics{./images/cyclohexanol.eps} & b & \num{110.6} & \num{160.8} \\ \hline

\end{longtable}

\begin{figure}[H]
\centering
	\caption{Ensayo de Lucas}
	\label{fig:Lucas}
	\includegraphics[width=10cm]{./images/Lucas_test.jpg} \\
\raggedright
\hspace{10pt} \small{FUENTE: Descargado de Internet bajo licencia \textit{Creative Commons}}
\end{figure}


\section{Discusión\label{sec:discusion}}

%La discusión debe dividirse en secciones de acuerdo a los experimentos realizados siguiendo una secuencia lógica. Todas las tablas y las figuras presentadas en la sección de resultados deben aparecer referenciadas. 

\subsection{Corrección de los puntos de ebullición}

Los puntos de ebullición fueron corregidos utilizando la ecuación de Sydney-Young\autocite{young:1902} (Ecuación \ref{eq:bpcorrection}). En donde \textit{t} es el punto de ebullición experimental, $C = \dfrac{dP}{dt} \times \dfrac{1}{T}$ es un valor aproximadamente constante que puede aproximarse a \num{0.00012} para líquidos polares y a \num{0.00010} para líquidos no polares y \textit{p} es la presión atmosférica del laboratorio expresada en mmHg.

\begin{equation}
\label{eq:bpcorrection}
\Delta T = C(760 - p)(273 + t)
\end{equation}

Se encontraron desviaciones respecto a los valores experimentales reportados en la literatura como se observa en la \ref{tab:boiling}. La diferencia observada se le puede atribuir al menos a tres razones; primero la presión atmosférica del laboratorio en el momento del experimento no se pudo medir así que se tomó un valor de \SI{680}{\mmHg}. Segundo el termómetro que se utilizó no fue calibrado previamente además que era un termómetro de inmersión completa y en el experimento solo estaba sumergido el bulbo, por lo que no existe certeza sobre la exactitud de la medida. Finalmente, el valor hay que recordar que los valores de la constante \textit{C} son aproximados y que en realidad esta constante es diferente para cada líquido.

\subsection{Ensayo de Lucas}

Como se observa en la \ref{fig:Lucas} en el tubo de la izquierda que contenía etanol no se observó reacción incluso después de 10 minutos, mientras que en el tubo de la derecha que contenía \textit{terc}-butanol la reacción ocurrió rápidamente y se observó la formación de dos fases. La turbidez que se observa en la capa inferior se debe a que el haluro de alquilo no se ha separado completamente y está parcialmente emulsionado. 

\section{Conclusiones\label{sec:conclusions}}

%Las conclusiones deben limitarse al resultado de los experimentos, deben ser claras y concisas.

\begin{itemize}
\item Primera conclusión
\item Segunda conclusión
\item \ldots etc.
\end{itemize}

\section{Referencias\label{sec:references}}

%Se deben consultar al menos cinco referencias bibliográficas relacionadas con la práctica

\printbibliography[heading=none]
		
\end{document}
